\chapter{Metodologia}\label{sec:metodos}

\begin{dica}[Dica de metodologia científica: o que escrever na Metodologia]
A metodologia descreve \emph{o que foi feito e como}, com detalhe suficiente para que
outra pessoa possa repetir ou avaliar o trabalho. Uma estrutura comum:
\begin{enumerate}
  \item \textbf{Classificação da pesquisa}: natureza (básica ou aplicada), abordagem
        (qualitativa, quantitativa ou mista), objetivos (exploratória, descritiva ou
        explicativa) e procedimentos (pesquisa bibliográfica, estudo de caso, pesquisa-ação,
        experimento, \emph{survey}, \emph{design science}, desenvolvimento de sistema).
        Justifique a escolha; não basta listar rótulos.
  \item \textbf{Etapas}: descreva a sequência de atividades, de preferência com um
        diagrama (\autoref{fig:etapas}) ligando cada etapa a um objetivo específico.
  \item \textbf{Instrumentos e materiais}: ferramentas, linguagens, versões, bases de
        dados, questionários e roteiros de entrevista (estes vão para apêndices).
  \item \textbf{Critérios de avaliação}: métricas, testes e como os resultados serão
        analisados (por exemplo, a escala \ac{sus} ou o \ac{tam}, tempo de resposta,
        cobertura de testes).
  \item \textbf{Aspectos éticos}: pesquisa com pessoas exige termo de consentimento e,
        em geral, avaliação do Comitê de Ética (Plataforma Brasil); trate dados pessoais
        conforme a \ac{lgpd}.
  \item \textbf{Reprodutibilidade e limitações}: publique o código em repositório com
        versão identificada, descreva o ambiente e discuta as ameaças à validade.
\end{enumerate}
Em trabalhos de Sistemas de Informação, detalhe também o processo de desenvolvimento
(ágil, iterativo etc.), o levantamento de requisitos (um quadro com requisitos funcionais
e não funcionais ajuda), a modelagem, a arquitetura, as tecnologias e a estratégia de testes.
\end{dica}

\section{Etapas da pesquisa}

Use figuras para esclarecer o texto. A \autoref{fig:etapas} resume as etapas da
pesquisa, e cada etapa corresponde a um objetivo específico.

\begin{figure}[htbp]
  \begin{ilustracao}{\caption{Etapas da metodologia}\label{fig:etapas}}{\elaboracaopropria}
    \begin{tikzpicture}[node distance=7mm, font=\small,
        etapa/.style={draw, rounded corners=2pt, minimum height=9mm, minimum width=24mm,
                      align=center, fill=black!5},
        seta/.style={-{Stealth[length=2.2mm]}, thick}]
      \node[etapa] (a) {Levantamento\\de requisitos};
      \node[etapa, right=of a] (b) {Modelagem};
      \node[etapa, right=of b] (c) {Implementação};
      \node[etapa, right=of c] (d) {Avaliação};
      \draw[seta] (a) -- (b);
      \draw[seta] (b) -- (c);
      \draw[seta] (c) -- (d);
    \end{tikzpicture}
  \end{ilustracao}
\end{figure}

\begin{dica}[Dica de uso da classe: figuras]
Na ilustração, o \emph{cabeçalho} (tipo, número e título) é centralizado, acima, e o
\emph{rodapé} (fonte, legenda e notas) fica abaixo, alinhado à margem esquerda da
própria ilustração, em espaço simples e tamanho 10 (manual UFV 4.9). O ambiente
\texttt{ilustracao}, do \texttt{abntex2-UFV}, mede o conteúdo e cuida disso. Use-o dentro
do \texttt{figure}, com os argumentos na ordem \emph{título} e \emph{rodapé}:
\texttt{\textbackslash begin\{ilustracao\}\allowbreak\{\textbackslash caption\{Título\}\allowbreak\textbackslash label\{rótulo\}\}\allowbreak\{\textbackslash fonte\{Fonte\}\allowbreak\textbackslash nota\{Nota\}\}},
com o \texttt{\textbackslash includegraphics} ou o diagrama dentro. Para a fonte de autoria
própria, use \texttt{\textbackslash elaboracaopropria}. Com
\texttt{\textbackslash caption[título curto]\{título longo\}}, a lista de ilustrações
recebe o nome curto. Prefira \ac{pdf} ou \ac{svg} convertido (vetoriais) para diagramas e
\ac{png} com pelo menos 300 dpi para capturas de tela, e sempre cite a figura no texto,
com \texttt{\textbackslash autoref}. Diagramas simples podem ser desenhados no próprio
\LaTeX\ com TikZ, como a \autoref{fig:etapas}.
\end{dica}

A \autoref{fig:logo} mostra a inserção de uma imagem; quando a imagem não é de autoria
própria, indique a fonte com \texttt{\textbackslash fonte} e inclua a obra nas referências.

\begin{figure}[htbp]
  \begin{ilustracao}{\caption[Exemplo de figura]{Exemplo de inserção de figura no \LaTeX}\label{fig:logo}}{\fonte{Universidade Federal de Viçosa (adaptado)}\nota{A imagem é um exemplo e deve ser substituída.}}
    \includegraphics[width=.7\linewidth]{Logo_Sistemas}
  \end{ilustracao}
\end{figure}

\section{Código-fonte}

Trechos de programa usam o ambiente \texttt{codigo}, que gera a ``Lista de códigos''
(\autoref{cod:exemplo}), como esta consulta \ac{sql} enviada a um \ac{sgbd} por meio de
uma \ac{api}. Programas extensos devem ir para apêndice ou repositório.

\begin{codigo}[htbp]
\begin{ilustracao}[\linewidth]{\caption{Exemplo de consulta SQL}\label{cod:exemplo}}{\elaboracaopropria}
\begin{lstlisting}[language=SQL, frame=tb]
SELECT nome, COUNT(*) AS total
FROM aluno
GROUP BY nome
ORDER BY total DESC;
\end{lstlisting}
\end{ilustracao}
\end{codigo}

\begin{dica}[Dica de uso da classe: código-fonte]
Escreva o código dentro de \texttt{lstlisting} (pacote \texttt{listings}) e informe a
linguagem em \texttt{[language=...]}. O ambiente \texttt{codigo}, definido pelo
\texttt{abntex2-UFV}, cria o rótulo ``Código N'' e a lista própria. Como a moldura do
\texttt{lstlisting} ocupa a largura da página, use
\texttt{\textbackslash begin\{ilustracao\}[\textbackslash linewidth]\{...\}\{...\}}: o argumento
opcional dá a largura mínima do título e do rodapé.
Para código com muitas linhas, \texttt{\textbackslash lstinputlisting\{arquivo.py\}}
lê o arquivo diretamente.
\end{dica}
